\documentclass[11pt,a4paper]{article}
\usepackage[spanish]{babel}
\usepackage[T1]{fontenc}
\usepackage[utf8]{inputenc}
\usepackage{lmodern}
\usepackage{geometry}
\usepackage{booktabs}
\usepackage{longtable}
\usepackage{tabularx}
\usepackage{array}
\usepackage{hyperref}
\usepackage{enumitem}
\usepackage{microtype}
\geometry{margin=2.5cm}
\hypersetup{colorlinks=true,urlcolor=blue,linkcolor=black}
\setlength{\parskip}{0.55em}
\setlength{\parindent}{0pt}
\title{\textbf{Impacto de la cuantización de modelos de lenguaje en sistemas RAG para consultas sobre documentación técnica industrial}\\[0.5em]\large Propuesta de proyecto final}
\author{Maestría en Robótica e Inteligencia Artificial (PRIA)\\UTEC --- Asignatura: Grandes Modelos de Lenguaje (LLM)}
\date{23 de noviembre de 2026}
\begin{document}
\maketitle
\section{Identificación del grupo y título tentativo}
\textbf{Integrantes:} Juan Pedro de León; \emph{[completar nombre del segundo integrante]}; \emph{[tercer integrante, si corresponde]}.

La asignatura exige grupos de dos a tres integrantes. La identificación se completará antes de entregar la propuesta.

\textbf{Título tentativo:} Impacto de la cuantización de modelos de lenguaje en sistemas RAG para consultas sobre documentación técnica industrial.

\textbf{Código del experimento:} \texttt{BENCH-QUANT-RAG-001}.

\section{Resumen de la propuesta}
La consulta de documentación técnica industrial mediante lenguaje natural puede facilitar el acceso a información sobre programación y automatización. Los grandes modelos de lenguaje (LLMs) ofrecen esta capacidad, pero su ejecución local está condicionada por la memoria disponible y los costos de inferencia. La cuantización reduce la precisión de representación de los pesos y los requisitos de memoria, aunque podría modificar la exactitud y la fidelidad de las respuestas.

Se propone un estudio experimental de seis representaciones del mismo modelo, \texttt{Qwen3-4B-Instruct-2507}: FP16, Q8\_0, Q6\_K, Q4\_K\_M, Q3\_K\_M y Q2\_K, integradas en un sistema de generación aumentada por recuperación (\emph{Retrieval-Augmented Generation}, RAG). Las consultas se responderán utilizando cinco fragmentos recuperados de manuales de Arduino y Siemens LOGO!. Se mantendrán constantes el corpus, el índice, los contextos recuperados y la configuración de inferencia para comparar las cuantizaciones.

El objetivo es cuantificar los compromisos entre exactitud factual, respaldo documental, variabilidad de respuestas, latencia y memoria de GPU, e identificar configuraciones que permitan ejecutar localmente consultas técnicas con recursos acotados. Se empleará un servidor con GPU AMD Radeon RX 9060 XT (16 GB), ROCm y \texttt{llama.cpp}.

El estudio se limita a una condición \emph{con RAG}. Por ello no buscará demostrar que RAG compense las pérdidas por cuantización frente a una condición sin recuperación.

\section{Relación con la asignatura}
El proyecto aplica conceptos de arquitectura Transformer y tokenización; representaciones vectoriales (\emph{embeddings}) para búsqueda semántica; ingeniería de prompts mediante instrucciones uniformes; y optimización de LLMs mediante cuantización. Integra asimismo RAG y evaluación de modelos en tareas concretas.

Como profundización, propone un diseño pareado de evaluación de calidad y eficiencia bajo distintas precisiones, con control de versiones, hashes, registro de recursos y separación estricta entre los datos de recuperación y las respuestas de referencia.

\section{Datos previstos}
Se utilizarán exclusivamente dos documentos PDF en el índice RAG: \texttt{Manual+Programacion+Arduino.pdf} (70 páginas) y \texttt{Logo\_s.pdf} (322 páginas, manual Siemens LOGO! edición 06/2003). El corpus ya fue extraído en 392 páginas y segmentado en 906 fragmentos de 900 caracteres con solapamiento de 150 caracteres. La recuperación utiliza \texttt{intfloat/multilingual-e5-small}, con vectores de 384 dimensiones e índice fijo.

Se dispone de 50 preguntas de evaluación: 20 de Arduino y 30 de LOGO!, con páginas PDF de referencia y respuestas esperadas. Estas referencias se mantienen fuera del índice y \textbf{no se incorporarán a los prompts}. Los documentos son históricos; los resultados se contrastarán prioritariamente con la versión exacta del manual utilizado y no se extrapolarán sin advertencia a equipos actuales.

La evaluación preliminar de recuperación Top-K=5 obtuvo, en ambas políticas evaluadas, Hit@5 de 0,96, MRR@5 de 0,8123 y recuperación media de páginas de referencia de 0,95. Las preguntas LOGO-012 y LOGO-024 no recuperaron la página de referencia entre los primeros cinco resultados; se conservarán como casos de recuperación insuficiente. La inferencia principal empleará la política global, sin filtro de dominio.

Los archivos de gran tamaño, los PDF, el banco con respuestas y las salidas completas permanecerán en el almacenamiento experimental local. GitHub contendrá el protocolo, los scripts y los manifiestos de procedencia que permitan verificarlos.

\section{Metodología tentativa}
\subsection{Arquitectura del sistema RAG y herramientas utilizadas}
El sistema RAG fue implementado mediante scripts Python independientes y verificables, sin depender de un framework de orquestación como LangChain o LlamaIndex ni de una base de datos vectorial externa. Esta decisión permite controlar y auditar la extracción, la segmentación, los embeddings, el ranking y el contexto recibido por el LLM.

\begin{center}
\small
\begin{tabularx}{\textwidth}{@{}p{0.25\textwidth}p{0.28\textwidth}X@{}}
\toprule
Componente & Herramientas & Función concreta\\
\midrule
Lectura y extracción & Python 3.12 y PyMuPDF & Extraer texto de cada página física del PDF, conservando su procedencia.\\
Segmentación & Python, algoritmo propio & Ventanas de hasta 900 caracteres con 150 de solapamiento, sin mezclar páginas.\\
Embeddings & SentenceTransformers y \texttt{multilingual-e5-small} & Transformar fragmentos y preguntas en vectores semánticos de 384 dimensiones.\\
Almacenamiento y búsqueda & NumPy, \texttt{embeddings.npy} & Matriz densa normalizada y ordenamiento por similitud coseno.\\
Evaluación de recuperación & Python, referencias de página & Calcular Hit@5, MRR@5 y cobertura de páginas del manual.\\
Generación & llama.cpp, GGUF y ROCm & Inferencia local del modelo Qwen3-4B en seis precisiones.\\
Trazabilidad & JSONL, JSON, CSV y SHA-256 & Congelar entradas, verificar integridad y registrar resultados.\\
\bottomrule
\end{tabularx}
\end{center}

\subsection{Preprocesamiento y extracción documental}
El corpus de entrada consta de dos manuales PDF: 70 páginas de Arduino y 322 de Siemens LOGO!, para un total de 392 páginas físicas. Mediante \texttt{prepare\_corpus.py} y la biblioteca \texttt{PyMuPDF}, se extrae el texto página por página con la opción \texttt{get\_text("text", sort=True)}. Se conserva el identificador del documento, el archivo original, el número de página PDF, el texto y su hash SHA-256. Se genera un inventario de fuentes (\texttt{source\_inventory.json}) y un registro de páginas (\texttt{pages.jsonl}).

Esta extracción utiliza el contenido textual del PDF; no se incorpora un proceso OCR. No se realizan operaciones de limpieza capaces de modificar silenciosamente el texto original. La extracción puede conservar imperfecciones de formato propias de la estructura PDF, una limitación que será considerada durante el análisis.

\subsection{Segmentación y conservación de procedencia}
El script \texttt{chunk\_corpus.py} aplica una segmentación determinista con ventanas máximas de 900 caracteres y 150 caracteres de solapamiento. Cuando es posible, ajusta el límite del fragmento a un espacio entre palabras. Cada fragmento pertenece a una única página: no se combinan textos de páginas distintas.

Se obtuvieron 906 fragmentos, identificados mediante un código que contiene documento, página y número de fragmento (por ejemplo, \texttt{logo-p0047-c001}). Cada registro conserva posiciones de caracteres, procedencia, texto y SHA-256. Los resultados se almacenan en \texttt{chunks.jsonl} y \texttt{chunk\_manifest.json}.

\subsection{Modelo de embeddings e índice vectorial}
Se emplea \texttt{intfloat/multilingual-e5-small}, disponible mediante la biblioteca \texttt{SentenceTransformers}, en la revisión fija \texttt{614241f622f53c4eeff9890bdc4f31cfecc418b3}. Se generan representaciones de 384 dimensiones para los 906 fragmentos. La implementación utiliza los prefijos específicos del modelo: \texttt{passage:} para documentos y \texttt{query:} para preguntas.

Los embeddings se convierten a \texttt{float32}, se normalizan mediante la norma L2 y se conservan en una matriz NumPy (\texttt{embeddings.npy}). No se utiliza FAISS, Chroma ni otra base vectorial: el índice corresponde a la matriz de vectores y sus metadatos de integridad. La búsqueda calcula la similitud coseno mediante el producto entre los vectores normalizados de la consulta y de los fragmentos, y ordena los resultados de mayor a menor puntuación.

\subsection{Recuperación, evaluación y congelamiento}
Para cada pregunta se seleccionan los cinco fragmentos mejor posicionados (\textbf{Top-K=5}). Se compararon dos estrategias: búsqueda global en ambos manuales y búsqueda restringida al dominio del documento. Ambas obtuvieron, en los 50 ítems, Hit@5 de 0,96, MRR@5 de 0,8123 y cobertura media de páginas de 0,95. Se eligió la política \emph{global} para la inferencia principal, evitando utilizar el dominio como información adicional para limitar la búsqueda.

El script \texttt{evaluate\_retrieval.py} separa la evaluación de la recuperación de los contextos que recibirán los modelos. Los archivos \texttt{retrieval\_evaluation.jsonl} contienen las métricas respecto de las páginas de referencia; \texttt{frozen\_contexts.jsonl} almacena exclusivamente preguntas y fragmentos recuperados, \textbf{sin respuestas de referencia ni páginas objetivo}. Las preguntas LOGO-012 y LOGO-024 no alcanzaron la página de referencia en el Top-K y se mantienen como casos de recuperación insuficiente.

El script \texttt{freeze\_rag\_prompts.py} comprueba los hashes y organiza los cinco fragmentos en un prompt reproducible, con citas de documento y página. El mismo texto y orden de evidencia se empleará con todas las cuantizaciones y repeticiones, sin recalcular embeddings ni modificar los documentos.

\subsection{Archivos y controles de reproducibilidad}
Los componentes y su procedencia se registran mediante identificadores, hashes SHA-256 y manifiestos JSON. Las principales salidas son:
\begin{itemize}[nosep]
    \item \texttt{source\_inventory.json} y \texttt{pages.jsonl}: fuentes y texto extraído.
    \item \texttt{chunks.jsonl} y \texttt{chunk\_manifest.json}: segmentación y parámetros.
    \item \texttt{embeddings.npy} y \texttt{index\_manifest.json}: matriz vectorial, revisión E5 y hashes.
    \item \texttt{frozen\_contexts.jsonl} y \texttt{retrieval\_manifest.json}: recuperación y su evaluación.
    \item Planillas CSV y manifiestos de modelos: tamaños, variantes y hashes de los GGUF.
\end{itemize}
Los documentos y archivos de evaluación originales se conservan fuera del repositorio de código, en el almacenamiento experimental local.

\subsection{Variantes del modelo}
Se utilizarán seis archivos GGUF de \texttt{Qwen3-4B-Instruct-2507} provenientes de la misma revisión fija del repositorio \texttt{unsloth/Qwen3-4B-Instruct-2507-GGUF}: \texttt{a06e946bb6b655725eafa393f4a9745d460374c9}. Se registrarán los SHA-256 individuales, tamaños, versión de llama.cpp, plantilla y parámetros de inferencia. Compartir repositorio no constituye por sí solo prueba de identidad del procedimiento de cuantización; su procedencia se declarará como limitación cuando corresponda.

\begin{center}
\begin{tabular}{ll}
\toprule
Variante & Función en el estudio\\
\midrule
FP16 & Referencia de mayor precisión\\
Q8\_0 & Cuantización de 8 bits\\
Q6\_K & Cuantización K de 6 bits\\
Q4\_K\_M & Cuantización K mixta de aproximadamente 4 bits\\
Q3\_K\_M & Cuantización K mixta de aproximadamente 3 bits\\
Q2\_K & Cuantización K de aproximadamente 2 bits\\
\bottomrule
\end{tabular}
\end{center}

\subsection{Diseño experimental}
Se consideran 50 preguntas, 6 variantes y 3 repeticiones por combinación, totalizando \textbf{900 respuestas}. Las variantes se ejecutarán individualmente. Se fijarán el prompt, la temperatura, el límite de generación, el contexto, la política de caché y, cuando corresponda, las semillas de muestreo. Cada ejecución conservará salida cruda, identificadores, tokens, métricas temporales, estado de finalización, errores y recursos de GPU. La primera prueba funcional de FP16 ya respondió una consulta sobre \texttt{setup()} y \texttt{loop()} mediante la API; no se considera una medición RAG formal.

\subsection{Evaluación y análisis}
La calidad se analizará en términos de exactitud frente a la respuesta de referencia, fidelidad de las afirmaciones al contexto recuperado, información no sustentada y variabilidad entre repeticiones. Se establecerá una rúbrica explícita de puntuación antes de analizar las 900 respuestas, para evitar adaptar el criterio a los resultados.

La eficiencia se analizará mediante latencia, tiempos de procesamiento y generación, tokens por segundo, tamaño del GGUF y consumo de VRAM. Para VRAM se distinguirán medición basal, uso tras carga y máximo durante inferencia, considerando que \texttt{rocm-smi} registra utilización de toda la GPU.

Las comparaciones serán pareadas por pregunta; las tres repeticiones no se tratarán como 150 preguntas independientes. Se presentarán tablas, gráficos y medidas de dispersión, explicitando posibles efectos del contenido recuperado, caché y variabilidad de generación.

\section{Viabilidad y resultados esperados}
Están disponibles los PDF, el corpus extraído y segmentado, el índice vectorial, los 50 ítems de evaluación y los seis GGUF. El servidor dispone de ROCm y llama.cpp. Las seis variantes GGUF superaron una prueba funcional de carga e inferencia (6/6, sin errores), aunque todavía no se han evaluado formalmente con el banco RAG. Se validó preliminarmente la API con FP16 y se observó una velocidad de 34,08 tokens/s en una consulta fuera del protocolo formal; esta cifra no constituye aún un resultado comparativo.

Se espera producir una caracterización reproducible del compromiso entre calidad y eficiencia para el modelo y el corpus seleccionados. Las conclusiones no se extenderán automáticamente a otros modelos, hardware, versiones de manuales o configuraciones RAG.

\section{Entregables}
El proyecto contempla un \textbf{informe técnico} con metodología, configuración, resultados, análisis, limitaciones y conclusiones, y una \textbf{presentación en línea de 20 minutos}, prevista para el \textbf{18 de diciembre de 2026}. El registro experimental incorporará manifiestos de hashes y planillas CSV generadas mediante scripts reproducibles.

\end{document}
